\newif\ifglobal

%%%%%%%%%%%%%%%%%%%%%%%
%%% Compile to view %%%
%%%%%%%%%%%%%%%%%%%%%%%

%%%%%%%%%%%%%%%%%%%%%%%%%%%%%%%%%%%%%%%%%%%%%%%%%%%
%%% Readme:                                     %%%
%%% https://www.overleaf.com/read/bwdjvfjksvjy/ %%%
%%%%%%%%%%%%%%%%%%%%%%%%%%%%%%%%%%%%%%%%%%%%%%%%%%%

% >>> M?: marks a module setting number or important notice

% >>> M4: Set {./relative-path} to external files for this module <<<
\def \modulefiles {./23-ECC}
% To add an external file, use:
% (Replace '---' with actual filename)
% '\input{\modulefiles/tables/---.tex}' for tables
% '\includegraphics{\modulefiles/figures/---}' for figures
% '\subfile{\modulefiles/---}' for register files


% >>> M1: This module is in global project? <<<
% 'YES' - uncomment; 'NO' - comment out
\globaltrue

\ifglobal
    \documentclass[main\_\_CN.tex]{subfiles}
   
\else
    % Font "FandolSong-Regular" used by ctex does not have GSUB table.
    % It causes warnings that do not affect compilation and can be ignored.
    \PassOptionsToPackage{quiet}{fontspec}

    \documentclass[openany, 10pt]{book}
   
    % >>> M2: In ./00-shared/preamble-trm-module, also find
    %         settings for visibility of todo notes, watermark <<< 
    \usepackage{preamble-trm-module}


    % Enable Chinese typesetting
    \usepackage{ctex}

    % >>> M3: Temporary labels in Chinese
    %         you can add your own labels there <<<
    \myexternaldocument{./00-shared/temp-labels-cn}

\fi %\ifglobal


\setcounter{secnumdepth}{4}


% >>> M5: If this module needs any updates in
% './00-shared/preamble-trm-module.sty'
% (include additional package, etc.), contact your module PM
% or create the file './00-shared/changelog.tex' make the
% necessary updates yourself and document them in this file <<<


\begin{document}


% Sets language into which pre-defined bits of text are translated
\selectlanguage{chinese}

% Do not delete this block
\ifglobal
% Add the following front matter if
% this module is outside of global project
\else
    \InputIfFileExists{readme}{}{}
    {\let\clearpage\relax \listoftodos}
    {\let\clearpage\relax \listofdonetodos}
    \tableofcontents
    %\listoftables
    %\listoffigures
\fi


%%%%%%%%%%%%%%%%%%%%%%%%%
%%% TRM Module Begins %%%
%%%%%%%%%%%%%%%%%%%%%%%%%

\hypertarget{ecc}{}
\chapter{ECC 加速器 (ECC)}
\label{mod:ecc}

\section{概述}
椭圆曲线密码学 (Elliptic Curve Cryptography) 是一种基于椭圆曲线数学的公开密钥加密演算法，其优势在于相对于 RSA 算法，使用较小长度的密钥就能够提供相当等级的加密安全性。

\chipname{} ECC 硬件加速器支持对于可选曲线的多种基础运算，用以实现对 ECC 基本运算、衍生算法（如 ECDSA 等算法）的加速。

\section{主要特性}
\chipname{} ECC 硬件加速器支持以下功能：

\begin{itemize}
    \item 支持两种可选 ECC 曲线，即 \href{https://nvlpubs.nist.gov/nistpubs/FIPS/NIST.FIPS.186-5.pdf}{FIPS 186-5} 中定义的 P-192 和 P-256
    \item 提供六种可选工作模式
    \item 提供计算完成的中断和中断控制
\end{itemize}

\section{专业名词定义}
本节介绍了 ECC 硬件加速器章节中会应用到的专业名词。

\subsection{ECC 背景知识}

\subsubsection{椭圆曲线与曲线上的点}

ECC 是一种基于大素数的有限域椭圆曲线的算法，这一椭圆曲线的数学表达式为：
\[
{y}^{2} = {x}^{3} + {ax} + b \ \textrm{mod}\ p
\]
其中，
\begin{itemize}
    \item $p$ 是素数，
    \item $a$ 和 $b$ 为两个小于 $p$ 的非负整数，
    \item $(x, y)$ 为满足该椭圆曲线方程的点。
\end{itemize}

\subsubsection{仿射坐标系与Jacobian坐标系}

一条椭圆曲线：
\begin{itemize}
    \item 在仿射坐标系下的表达式为：
    \[
    {y}^{2} = {x}^{3} + {ax} + b \ \textrm{mod}\ p
    \]
    \item 在 Jacobian 坐标系下的表达式为：
    \[
    {Y}^{2} = {X}^{3} + {aXZ}^{4} + {bZ}^{6} \ \textrm{mod}\ p
    \]
\end{itemize}

一个曲线上的点在仿射坐标系下的表示 $(x, y)$ 与在 Jacobian 坐标系下的表示 $(X, Y, Z)$ 有如下转换关系：

\begin{itemize}
    \item 从 Jacobian 坐标系到仿射坐标系的转换：
    \[
    x = X / {Z}^2 \ \textrm{mod}\ p
    \]
    \[
    y = Y / {Z}^3 \ \textrm{mod}\ p
    \]
    \item 从仿射坐标系到 Jacobian 坐标系的转换：
    \[ X = x \]
    \[ Y = y \]
    \[ Z = 1\]
\end{itemize}


\subsection{\chipname{} ECC 相关定义}

\subsubsection{内存块}
ECC 硬件加速器的内存块用于存储 ECC 运算中所用到的输入数据和输出数据。
\begin{table}[H]
    \centering
    \caption{ECC 硬件加速器内存块} \label{tab:ecc-memory}
    \begin{threeparttable}
    \begin{tabular}{ | l | l | l | l | l | l | }
    \hline
        \rowcolor{lightgray}
        \textbf{内存块} & \textbf{大小 (byte)}  &  \textbf{起始地址}\tnote{*} &  \textbf{结束地址} \tnote{*} & \textbf{访问属性}   \\ \hline
        ECC\_MULT\_Mem\_k    & 32  & 0x100 & 0x11F & R/W \\ \hline
        ECC\_MULT\_Mem\_Px   & 32  & 0x120 & 0x13F & R/W \\ \hline
        ECC\_MULT\_Mem\_Py   & 32  & 0x140 & 0x15F & R/W \\ \hline
    \end{tabular}
    \begin{tablenotes}
        \item[*] 采用相对于 ECC 加速器基地址的偏移量。详见章节 \ref{mod:sysmem} \textit{\nameref{mod:sysmem}} 中的表 \ref{tab:sysmem-base-address}。
    \end{tablenotes}
    \end{threeparttable}
\end{table}

\subsubsection{数据与数据块}

在 ECC 硬件加速器模块中会用到的数据位宽均为 256 位，假设一个数据为 $D[255:0]$，则其可以被划分为八个 32 位的数据块 $D[n][31:0](n=0,1,\cdots,7)$，序号数低的数据块对应二进制低位，即：
\[
D[255:0] = {D[7][31:0],D[6][31:0],D[5][31:0],
D[4][31:0],D[3][31:0],D[2][31:0],D[1][31:0],D[0][31:0]}
\]

\subsubsection{数据存储}

数据存储即将一个数据存储进一个内存块的操作，也可以说该数据为 ECC 算法的输入数据。具体来说，将数据写入一个 ECC 内存块相当于将该数据 $D[n][31:0] (n=0, 1, \cdots, 7)$ 依次写入“该内存块起始地址 + $
4 \times n$”：

\begin{itemize}
    \item 写入 $D[0]$ 至“起始地址”
    \item 写入 $D[1]$ 至“起始地址 + 4”
    \item $\cdots$
    \item 写入 $D[7]$ 至“起始地址 + 28”
\end{itemize}
\begin{tiplisting}
在 192 位模式下，进行数据存储操作时，需要在数据的高位补 0，保证存储的数据为 256 位。
\end{tiplisting}

\subsubsection{数据读取}

数据读取即将一个数据从一个内存块读出的操作，也可以说该数据为 ECC 算法的输出数据。具体来说，从一个 ECC 内存块读数据相当于从“该内存块起始地址 + $
4 \times n$”依次读出 $D[n][31:0] (n=0, 1, \cdots, 7)$：
\begin{itemize}
    \item 从“起始地址”读出 $D[0]$
    \item 从“起始地址 + 4”读出 $D[1]$
    \item $\cdots$
    \item 从“起始地址 + 28”读出 $D[7]$
\end{itemize}

\begin{tiplisting}
在 192 位模式下，进行数据读取操作时，只需要读取低 192 位（即 6 个数据块）的数据。
\end{tiplisting}

\subsubsection{标准运算与 Jacobian 运算}

\chipname{} ECC 硬件加速器中，所有标准运算（包括标准点验证和标准点乘）的输入数据以及输出数据中的点均在仿射坐标系中；相对应地，所有 Jacobian 运算（包括 Jacobian 点验证和 Jacobian 点乘）的输入数据以及输出数据中的点均在 Jacobian 坐标系中。

\section{功能描述}
\label{sec:ecc-mode-config}
\subsection{密钥长度模式}

\chipname{} ECC 硬件加速器共支持 2 种密钥长度模式，每种密钥长度模式唯一对应一条椭圆曲线。用户通过配置 \hyperref[fielddesc:ECCMULTKEYLENGTH]{ECC\_MULT\_KEY\_LENGTH} 可选定密钥长度模式，其与椭圆曲线的对应关系如表 \ref{tab:ecc-key-length} 所示。
\begin{table}[H]
    \centering
    \caption{ECC加速器密钥长度模式控制}
    \label{tab:ecc-key-length}
    \begin{threeparttable}
    \begin{tabular}{|p{4.2cm}|p{3.6cm}|}
    \hline
        \rowcolor{lightgray}
        \textbf{\hyperref[fielddesc:ECCMULTKEYLENGTH]{ECC\_MULT\_KEY\_LENGTH}} & \textbf{对应椭圆曲线\tnote{*}} \\ \hline
        1'b0 & FIPS P-192  \\\hline
        1'b1 & FIPS P-256  \\\hline
    \end{tabular}
    \begin{tablenotes}
        \item[*] FIPS P-192/P-256 的曲线定义，请参考 \href{https://nvlpubs.nist.gov/nistpubs/FIPS/NIST.FIPS.186-5.pdf}{FIPS 186-5}。

    \end{tablenotes}
    \end{threeparttable}
\end{table}

\subsection{工作模式}

\chipname{} ECC 硬件加速器共有六种工作模式，每种工作模式进行基于选定曲线的不同操作。用户通过配置寄存器 \hyperref[fielddesc:ECCMULTWORKMODE]{ECC\_MULT\_WORK\_MODE} 可选定工作模式，其与工作模式的对应关系如表 \ref{tab:ecc-work-mode} 所示。

\begin{table}[H]
    \centering
    \caption{ECC 硬件加速器工作模式控制}
    \label{tab:ecc-work-mode}
    \begin{tabular}{|l|l|l|l|}
    \hline
        \rowcolor{lightgray}
        \textbf{\hyperref[fielddesc:ECCMULTWORKMODE]{ECC\_MULT\_WORK\_MODE}} & \textbf{对应模式} & \textbf{\hyperref[fielddesc:ECCMULTWORKMODE]{ECC\_MULT\_WORK\_MODE}} & \textbf{对应模式} \\ \hline
        3'd0 & 标准点乘                & 3'd4 & Jacobian 点乘            \\\hline
        3'd1 & \textbf{保留项，不可用} & 3'd5 & \textbf{保留项，不可用}  \\\hline
        3'd2 & 标准点验证              & 3'd6 & Jacobian 点验证          \\\hline
        3'd3 & 标准点验证 + 标准点乘     & 3'd7 & 标准点验证 + Jacobian 点乘  \\\hline
    \end{tabular}
\end{table}

每种工作模式的具体计算和输入/输出数据，请参照下述子章节。
\subsubsection{标准点乘模式}

该模式计算如下公式：
\[
(Q_x,Q_y) = k \cdot (P_x,P_y)
\]
其中，
\begin{itemize}
    \item 输入数据：$P_x$，$P_y$，$k$对应的内存块为 \hyperref[tab:ecc-memory]{ECC\_MULT\_Mem\_Px}, \hyperref[tab:ecc-memory]{ECC\_MULT\_Mem\_Py} 和 \hyperref[tab:ecc-memory]{ECC\_MULT\_Mem\_k}。
    \item 输出数据：$Q_x$, $Q_y$ 对应的内存块为 \hyperref[tab:ecc-memory]{ECC\_MULT\_Mem\_Px} 和 \hyperref[tab:ecc-memory]{ECC\_MULT\_Mem\_Py}。
\end{itemize}



\subsubsection{标准点验证模式}

该模式用于计算点 $(P_x, P_y)$ 是否在选定的椭圆曲线上。
其中，
\begin{itemize}
    \item 输入数据：$P_x$，$P_y$ 对应的内存块为 \hyperref[tab:ecc-memory]{ECC\_MULT\_Mem\_Px} 和 \hyperref[tab:ecc-memory]{ECC\_MULT\_Mem\_Py}。
    \item 输出数据：点验证的结果存储在寄存器 \hyperref[fielddesc:ECCMULTVERIFICATIONRESULT]{ECC\_MULT\_VERIFICATION\_RESULT} 中。
\end{itemize}

\subsubsection{标准点验证+标准点乘模式}
该模式先计算点 $(P_x,P_y)$ 是否在选定的椭圆曲线上，如果其在选定的椭圆曲线上，则继续计算如下公式：
\[
(Q_x,Q_y) = k \cdot (P_x,P_y)
\]
其中，
\begin{itemize}
    \item 输入数据：$P_x$，$P_y$，$k$对应的内存块为 \hyperref[tab:ecc-memory]{ECC\_MULT\_Mem\_Px}, \hyperref[tab:ecc-memory]{ECC\_MULT\_Mem\_Py} 和 \hyperref[tab:ecc-memory]{ECC\_MULT\_Mem\_k}。
    \item 输出数据：
    \begin{itemize}
        \item 点验证的结果存储在寄存器 \hyperref[fielddesc:ECCMULTVERIFICATIONRESULT]{ECC\_MULT\_VERIFICATION\_RESULT} 中。
        \item $Q_x$、$Q_y$对应的内存块为 \hyperref[tab:ecc-memory]{ECC\_MULT\_Mem\_Px} 和 \hyperref[tab:ecc-memory]{ECC\_MULT\_Mem\_Py}。
    \end{itemize}
\end{itemize}

\subsubsection{Jacobian 点乘模式}
该模式计算如下公式： 
\[
(Q_x,Q_y,Q_z) = k \cdot (P_x,P_y,1)
\]
其中，
\begin{itemize}
    \item $(Q_x,Q_y,Q_z)$ 为 Jacobian 表示的曲线上的点。
    \item 输入点P的Jacobian表示中添加的1为硬件默认补全，不需要输入。
    \item 输入数据：$P_x$，$P_y$ 和 $k$对应的内存块为 \hyperref[tab:ecc-memory]{ECC\_MULT\_Mem\_Px}, \hyperref[tab:ecc-memory]{ECC\_MULT\_Mem\_Py} 和 \hyperref[tab:ecc-memory]{ECC\_MULT\_Mem\_k}。
    \item 输出数据：$Q_x$、$Q_y$ 和 $Q_z$对应的内存块为 \hyperref[tab:ecc-memory]{ECC\_MULT\_Mem\_Px}, \hyperref[tab:ecc-memory]{ECC\_MULT\_Mem\_Py} 和 \hyperref[tab:ecc-memory]{ECC\_MULT\_Mem\_k}。
\end{itemize}
\subsubsection{Jacobian 点验证模式}
该模式用于计算点$(Q_x,Q_y,Q_z)$是否在选定的椭圆曲线上。
其中，
\begin{itemize}
    \item $(Q_x,Q_y,Q_z)$ 为Jacobian表示的点。
    \item 输入数据：$Q_x$、$Q_y$ 和 $Q_z$对应的内存块为 \hyperref[tab:ecc-memory]{ECC\_MULT\_Mem\_Px}, \hyperref[tab:ecc-memory]{ECC\_MULT\_Mem\_Py} 和 \hyperref[tab:ecc-memory]{ECC\_MULT\_Mem\_k}。
    \item 输出数据：点验证的结果存储在寄存器 \hyperref[fielddesc:ECCMULTVERIFICATIONRESULT]{ECC\_MULT\_VERIFICATION\_RESULT} 中。
\end{itemize}

\subsubsection{标准点验证 + Jacobian 点乘模式}

该模式先计算点 $(P_x,P_y)$ 是否在选定的椭圆曲线上，如果其在选定的椭圆曲线上，则继续计算如下公式：
\[
(Q_x,Q_y,Q_z) = k \cdot (P_x,P_y,1)
\]
其中
\begin{itemize}
    \item $(Q_x,Q_y,Q_z)$ 为 Jacobian 表示的曲线上的点。
    \item 输入点 $P$ 的 Jacobian 表示中添加的 1 为硬件默认补全，不需要输入。
    \item 输入数据：$P_x$，$P_y$ 和 $k$对应的内存块为 \hyperref[tab:ecc-memory]{ECC\_MULT\_Mem\_Px}, \hyperref[tab:ecc-memory]{ECC\_MULT\_Mem\_Py} 和 \hyperref[tab:ecc-memory]{ECC\_MULT\_Mem\_k}。
    \item 输出数据：
    \begin{itemize}
        \item 点验证的结果存储在寄存器 \hyperref[fielddesc:ECCMULTVERIFICATIONRESULT]{ECC\_MULT\_VERIFICATION\_RESULT} 中。
        \item $Q_x$、$Q_y$ 和 $Q_z$ 对应的内存块为 \hyperref[tab:ecc-memory]{ECC\_MULT\_Mem\_Px}，\hyperref[tab:ecc-memory]{ECC\_MULT\_Mem\_Py} 和 \hyperref[tab:ecc-memory]{ECC\_MULT\_Mem\_k}。
    \end{itemize}
\end{itemize}

\section{时钟与复位} 
\chipname{} ECC 硬件加速器模块仅有一个模块时钟 CRYPTO\_ECC\_CLK 和一个模块复位 CRYPTO\_ECC\_RST。在使用 ECC硬件加速器之前，需要开启 ECC 时钟，关闭 ECC 复位。如何配置 ECC 时钟和复位，请参考章节 \ref{mod:resclk} \textit{\nameref{mod:resclk}}。

\section{中断}
\label{sec:ecc-interrupt}

\chipname{} ECC 硬件加速器可产生一个中断信号 ECC\_INTR，并将其发送给\hyperref[mod:intmtrx]{中断矩阵}。

\begin{tiplisting}
每个中断信号均由其包含的所有中断源的状态位共同产生，即任意一个其包含的中断源的状态位触发，该中断信号就会被触发。
\end{tiplisting}

ECC 硬件加速器的中断信号 ECC\_INTR 仅包含中断源 \label{int:ECCMULTCALCDONEINT}ECC\_MULT\_CALC\_DONE\_INT，ECC 硬件加速器运算完成即触发该中断源。中断源 ECC\_MULT\_CALC\_DONE\_INT 由以下寄存器控制：

\begin{itemize}
    \item \hyperref[fielddesc:ECCMULTCALCDONEINTRAW]{ECC\_MULT\_CALC\_DONE\_INT\_RAW}：ECC 硬件加速器运算完成时置 1。
    \item \hyperref[fielddesc:ECCMULTCALCDONEINTST]{ECC\_MULT\_CALC\_DONE\_INT\_ST}：反映 ECC 硬件加速器运算完成中断的状态，通过用 \hyperref[fielddesc:ECCMULTCALCDONEINTENA]{ECC\_MULT\_CALC\_DONE\_INT\_ENA} 使能/屏蔽 \hyperref[fielddesc:ECCMULTCALCDONEINTRAW]{ECC\_MULT\_CALC\_DONE\_INT\_RAW} 位来生成。
    \item \hyperref[fielddesc:ECCMULTCALCDONEINTENA]{ECC\_MULT\_CALC\_DONE\_INT\_ENA}：用于使能或屏蔽 ECC 硬件加速器运算完成中断状态位。
    \item \hyperref[fielddesc:ECCMULTCALCDONEINTCLR]{ECC\_MULT\_CALC\_DONE\_INT\_CLR}：置 1 此位清除 ECC 硬件加速器运算完成中断，对应的 \hyperref[fielddesc:ECCMULTCALCDONEINTRAW]{ECC\_MULT\_CALC\_DONE\_INT\_RAW} 和 \hyperref[fielddesc:ECCMULTCALCDONEINTST]{ECC\_MULT\_CALC\_DONE\_INT\_ST} 位会清零。
\end{itemize}

\section{软件配置流程}
软件配置 ECC 硬件加速器的流程如下：
\begin{enumerate}
    \item 配置 ECC 硬件加速器的时钟与复位。
    \item 根据 \ref{sec:ecc-mode-config} 小节的描述，按照需求配置ECC加速器密钥长度模式和工作模式。
    \item 根据 \ref{sec:ecc-interrupt} 小节的描述，使能 ECC\_MULT\_CALC\_DONE\_INT 中断。
    \item 置位寄存器 \hyperref[fielddesc:ECCMULTSTART]{ECC\_MULT\_START} 以启动 ECC 运算。
    \item 等待 ECC\_MULT\_CALC\_DONE\_INT 中断的产生，即 ECC 运算结束。
    \item 根据 \ref{sec:ecc-mode-config} 小节的描述，查看运算结果。
\end{enumerate}

% >>> If your module has no registers, delete sample text
%       for Sections Register Summary and Registers below <<<

%%%%%%%%%%%%%%%%%%%%%%%%
%%% Register Summary %%%
%%%%%%%%%%%%%%%%%%%%%%%%

% >>> Update the sample text below and insert
%     the info on register summary into the subfile <<<

\clearpage
\hypertarget{ecc-reg-summ}{}
\section{寄存器列表}
\label{sec:ecc-reg-summ}

本小节的所有地址均为相对于 ECC 加速器基地址的地址偏移量（相对地址），具体基地址请见章节 \ref{mod:sysmem} \textit{\nameref{mod:sysmem}} 中的表 \ref{tab:sysmem-base-address}。

请查看章节 \hyperref[glossary-access-types]{\textit{寄存器的访问类型}}，了解“访问”列缩写的含义。

\subfile{\modulefiles/23-ECC-reg-summ__CN}


%%%%%%%%%%%%%%%%%
%%% Registers %%%
%%%%%%%%%%%%%%%%%

% >>> Update the sample text below and insert
%      the info on registers into the subfile <<<

\clearpage
\section{寄存器}

本小节的所有地址均为相对于 ECC 加速器基地址的地址偏移量（相对地址），具体基地址请见章节 \ref{mod:sysmem} \textit{\nameref{mod:sysmem}} 中的表 \ref{tab:sysmem-base-address}。

\subfile{\modulefiles/23-ECC-reg__CN}


\end{document}
